\section{FAIRSEQ2}

FAIRSEQ2 is an open-source library of sequence modeling components that provides researchers and developers with building blocks for machine translation, language modeling, and other sequence generation tasks, specifically around text and audio data format. FAIRSEQ2 is distributed with an MIT license and is available on GitHub at \href{https://github.com/pytorch/fairseq2}{https://github.com/pytorch/fairseq2}.

FAIRSEQ2 features: 
(i) state-of-the-art implementations of transformers and their components (transformer layers, embedding layers, layernorms, attention blocks, etc.);% 
(ii) \texttt{fairseq2.data} --  a scalable pipeline API that enables text and audio data pre-processing, transformation, shuffling, and batching in a streaming manner, allowing training over multi-terabyte datasets without explicit data preparation steps or data loading timeouts; (iii) core building components for efficient model training (optimizers, LR schedulers, loss implementations); (iv) sequence generators for optimized inference with incremental beam search. %

Following the spirit of its predecessor FAIRSEQ~\citep{ott2019fairseq}, FARSEQ2 was built with extensibility in mind. The library-like structure of the code enables effortless component drop-ins, including those initially written in FAIRSEQ. 
We expect continuous populating of the library with new components by us and by the open-source community in the following years.
%

Another guiding principle for FAIRSEQ2 is a clear separation of core and experimental code. 
The original FAIRSEQ has become a hub for numerous research ideas. 
Often they were added in the form of if-else statements mixed with the core functionality. 
Over time, the number of such if-else statements and associated command line options has grown, with each option poorly supported and often subtly incompatible with other options. 
To prevent this scenario, in FAIRSEQ2, all basic components are designed with the “dependency inversion” principle, making it possible to compose them easily. 
Existing model architectures can be modified with just a few lines of code without requiring copy/pasting large amounts of code. All plug-ins and modifications exist as separate components, not interfering with the parent blocks and not hindering access to them for other users. Larger efforts (like \unity or \sonar described in this paper) are moved into separate repositories and use FAIRSEQ2 as a dependency.

We acknowledge the wide range of training and execution environments for Deep Learning models that exist today (from a single-container training via on-demand Cloud Computing Services to huge LLMs training jobs running on exaFLOPS supercomputers with tens of thousands GPUs; from very limited inference capabilities of edge devices to the power of accelerated inference on ASICs). To meet the diverse expectations of these environments, FAIRSEQ2 has shifted from the idea of a self-contained single-stop for all training, evaluation, and inference pipelines towards a set of independent components that can be used and extended outside of FAIRSEQ2. We put an emphasis on compatibility with the existing alternatives in PyTorch and other Deep Learning frameworks, following common API conventions and inheriting from the same base classes. That guarantees effortless drop-in replacement of components from different origins. The user is offered a wide range of usage scenarios: from implementing a complete pipeline using FAIRSEQ2 to fusing multiple Deep Learning frameworks in their project, or even picking a single block like the efficient implementation of an optimizer.
